\chapter{Introduction}

Enterprise telecom integration teams repeatedly build API connectors that map external REST contracts to internal EJB-based services. The lifecycle is predictable but fragmented: engineers parse OpenAPI/Swagger files, inspect compiled JAR metadata, map request and response fields, prepare YAML configurations, and run generation scripts. In many teams, these steps are handled with separate tools and ad-hoc notes, which increases onboarding time, introduces inconsistency, and raises maintenance cost.

This thesis studies the API Development Tool built during an internship at Amdocs. The platform combines a React 18 frontend~\cite{react_docs} and a Node.js/Express backend~\cite{express_docs} to support specification upload, JAR bytecode introspection, interactive mapping, YAML generation, and connector module generation with Socket.IO~\cite{socketio_docs} progress streaming. Its main contribution is workflow unification: previously disconnected activities become a single, traceable process.

\section{Motivation}

The original connector workflow relied heavily on manual interpretation and undocumented decisions. Engineers used text editors for API contracts, Java tools for class inspection, spreadsheets for mappings, and command-line scripts for generation. This approach consumed time and produced variation in naming, structure, and mapping quality.

Because downstream generation depends on strict YAML structure, mapping errors propagated into generated code and Maven artifacts. The process also lacked a central history of why mappings were chosen, making updates and debugging slow. These issues motivated a guided platform that enforces consistency and records decisions.

Table~\ref{tab:manual_vs_automated} summarises the transition from manual to tool-assisted development.

\begin{table}[H]
\centering
\caption{Comparison of manual and tool-assisted connector development.}
\label{tab:manual_vs_automated}
\small
\begin{tabular}{|p{3.2cm}|p{4.5cm}|p{4.5cm}|}
\hline
\textbf{Aspect} & \textbf{Manual Workflow} & \textbf{Tool-Assisted Workflow} \\
\hline
Specification analysis & Text editors, manual reading & Automated parsing with reference resolution and schema flattening \\
\hline
Service discovery & JAR decompilation in Java IDEs & Browser-based bytecode introspection without JVM \\
\hline
Field mapping & Spreadsheets, ad-hoc notes & Interactive dropdowns with source inference \\
\hline
Configuration authoring & Hand-written YAML files & Automated YAML generation from mapping state \\
\hline
Code generation & Command-line script invocation & One-click generation with real-time streaming \\
\hline
Traceability & Undocumented, scattered files & Full audit trail in database \\
\hline
Consistency & Varies by engineer & Enforced by structured workflow \\
\hline
\end{tabular}
\end{table}

\section{Problem Context}

The platform operates in an existing enterprise environment, not a greenfield system. It must align with a mature digital connector module tree, Maven-based build plugins, compiled client-kit JAR artifacts, and mixed Windows/Linux deployment conditions, including air-gapped production networks.

Accordingly, generated outputs must match fixed directory conventions, such as Swagger files under \texttt{src/\allowbreak main/\allowbreak resources/\allowbreak swagger/} and configs under \texttt{src/\allowbreak main/\allowbreak resources/\allowbreak configs/}. The key challenge is end-to-end continuity: connecting ingestion, introspection, mapping, generation, and synchronisation without breaking enterprise constraints.

\section{Objectives}

This thesis provides a code-level and architecture-level analysis of the platform. It traces how inputs move through frontend, backend, filesystem, and generation layers from specification upload to deployable output.

It also evaluates runtime and deployment behaviour, including Tomcat-hosted frontend delivery, PM2-managed backend operation, certificate handling, and filesystem-centric persistence. The study reports evidence-backed findings across security, reliability, portability, and maintainability, and proposes prioritised improvements.

\section{Method of Study}

The methodology is repository-driven technical analysis. Backend evaluation covers route logic, persistence design, synchronisation routines, and generation orchestration. Frontend evaluation covers workflow components, configuration editing, and API communication.

The Swagger parser, JAR parser, and configuration generator are examined in depth: reference resolution and schema flattening, bytecode and signature parsing, EJB detection, field-path extraction, source inference, and multi-EJB mapping assembly. Conclusions are grounded in observed implementation behaviour.

\section{Report Organisation}

The remaining chapters are organised as follows:

\begin{enumerate}
    \setcounter{enumi}{1}
    \item \textbf{Chapter 2: Purpose and Scope} defines project intent, functional boundaries, and expected outcomes.
    \item \textbf{Chapter 3: Related Works} positions the platform against API-first tooling, code generators, and workflow automation approaches.
    \item \textbf{Chapter 4: System Architecture and Design} details frontend, backend, utilities, synchronisation, generation, and deployment design.
    \item \textbf{Chapter 5: Results and Analysis} presents evidence-based evaluation across workflow coverage, architecture, reliability, security, and portability.
    \item \textbf{Chapter 6: Conclusion and Future Work} summarises contributions and outlines a phased enhancement roadmap.
\end{enumerate}
